\documentclass[a4paper,11pt]{article}
\usepackage[utf8]{inputenc}
\usepackage[T1]{fontenc}
\usepackage[french]{babel}
\usepackage{geometry}
\geometry{left=2cm,right=2cm,top=1.5cm,bottom=2cm}
\usepackage{amsmath,amssymb}
\usepackage{enumitem}
\usepackage{array}
\usepackage{booktabs}
\usepackage{xcolor}
\usepackage{tcolorbox}
\usepackage{fancyhdr}
\usepackage{tikz}
\usepackage{lastpage}

% Couleurs
\definecolor{bleupro}{HTML}{0969da}
\definecolor{orange}{HTML}{d97706}
\definecolor{vertpro}{HTML}{2f855a}
\definecolor{rouge}{HTML}{c53030}
\definecolor{gris}{HTML}{718096}

% En-tête
\pagestyle{fancy}
\fancyhf{}
\lhead{\small CCF Physique-Chimie — Terminale ICCER}
\rhead{\small Fluides \& Courant alternatif}
\cfoot{\thepage/\pageref{LastPage}}
\renewcommand{\headrulewidth}{0.5pt}

% Boîtes
\tcbuselibrary{skins,breakable}

\newtcolorbox{formulairebox}{
  colback=yellow!5,
  colframe=orange,
  title={\textbf{FORMULAIRE — Rappels de cours}},
  fonttitle=\large\bfseries,
  breakable
}

\newtcolorbox{contextebox}{
  colback=blue!3,
  colframe=bleupro,
  leftrule=4pt,
  rightrule=0pt,
  toprule=0pt,
  bottomrule=0pt,
  boxsep=4pt,
  left=8pt
}

\newtcolorbox{reponsebox}[1][4cm]{
  colback=white,
  colframe=gris!40,
  boxrule=0.5pt,
  arc=3pt,
  height=#1,
  valign=top,
  before skip=4pt,
  after skip=8pt
}

\newcommand{\pts}[1]{\hfill\textcolor{gris}{\small\textbf{#1}}}
\newcommand{\comp}[1]{\colorbox{blue!10}{\scriptsize\textbf{#1}}}

\setlength{\parindent}{0pt}
\setlength{\parskip}{6pt}

\begin{document}

% ===== PAGE DE GARDE =====
\begin{center}
{\LARGE\bfseries\color{bleupro} Contrôle en Cours de Formation}\\[8pt]
{\large Physique-Chimie — Terminale Bac Pro ICCER — Groupement 1}\\[12pt]
{\Large\bfseries Débit de fluide \& Puissance en régime sinusoïdal}\\[20pt]

\begin{tabular}{|p{6cm}|p{6cm}|}
\hline
\textbf{Durée :} 45 minutes & \textbf{Barème :} /10 points \\
\hline
\textbf{Documents :} Calculatrice autorisée & \textbf{Formulaire :} ci-dessous \\
\hline
\textbf{Nom :} \dotfill & \textbf{Prénom :} \dotfill \\
\hline
\textbf{Classe :} \dotfill & \textbf{Date :} \dotfill \\
\hline
\end{tabular}
\end{center}

\vspace{10pt}

% ===== FORMULAIRE =====
\begin{formulairebox}

\subsection*{Débit et transport de fluide}

\begin{tabular}{>{\bfseries}p{4.5cm} p{5.5cm} p{4cm}}
\toprule
\textbf{Grandeur} & \textbf{Formule} & \textbf{Unités} \\
\midrule
Débit volumique & $Q_v = S \times v$ & m³/s, m², m/s \\[4pt]
Débit massique & $Q_m = \rho \times Q_v$ & kg/s \\[4pt]
Équation de continuité & $S_1 \times v_1 = S_2 \times v_2$ & \\[4pt]
Section d'un tuyau & $S = \pi \times \left(\dfrac{d}{2}\right)^2$ & m² \\[6pt]
Conversions & 1 m³/h $= \dfrac{1}{3\,600}$ m³/s & \\[4pt]
\bottomrule
\end{tabular}

\smallskip
$\rho_{\text{eau}} = 1\,000$ kg/m³ \quad Vitesse recommandée dans les canalisations : entre 0,5 et 1,5 m/s

\subsection*{Puissance en régime sinusoïdal}

\begin{tabular}{>{\bfseries}p{4.5cm} p{5.5cm} p{4cm}}
\toprule
\textbf{Grandeur} & \textbf{Formule} & \textbf{Unités} \\
\midrule
Puissance active & $P = U \times I \times \cos\varphi$ & W \\[4pt]
Puissance apparente & $S = U \times I$ & VA \\[4pt]
Puissance réactive & $Q = U \times I \times \sin\varphi$ & VAR \\[4pt]
Facteur de puissance & $\cos\varphi = \dfrac{P}{S}$ & (sans unité) \\[6pt]
Énergie & $E = P \times t$ & J (ou kWh) \\[4pt]
\bottomrule
\end{tabular}

\smallskip
\textbf{Triangle des puissances :} $S^2 = P^2 + Q^2$. Un $\cos\varphi$ proche de 1 est souhaitable (moins de pertes réactives).

\end{formulairebox}

\newpage

% ================================================================
\section*{Exercice 1 — Dimensionnement d'un réseau de chauffage \pts{5 points}}
% ================================================================

\begin{contextebox}
Un installateur thermique dimensionne le réseau d'eau chaude d'un bâtiment. Le circulateur de la chaudière fournit un débit volumique $Q_v = 1{,}8$ m³/h. La canalisation principale a un diamètre $d_1 = 40$ mm. Elle se divise ensuite en deux branches identiques de diamètre $d_2 = 25$ mm chacune (le débit se répartit à parts égales).

La vitesse dans les canalisations doit rester entre 0,5 et 1,5 m/s pour limiter le bruit et l'usure.
\end{contextebox}

% Schéma
\begin{center}
\begin{tikzpicture}[scale=0.8, every node/.style={font=\small}]
  % Chaudière
  \draw[fill=red!10, rounded corners] (0,0.5) rectangle (2,2.5);
  \node[align=center] at (1,1.5) {\textbf{Chaudière}\\$Q_v = 1{,}8$ m³/h};
  % Canalisation principale
  \draw[very thick, bleupro] (2,1.5) -- (5,1.5);
  \node[above, bleupro] at (3.5,1.5) {$d_1 = 40$ mm};
  % Division
  \filldraw[bleupro] (5,1.5) circle (3pt);
  % Branche haute
  \draw[thick, vertpro] (5,1.5) -- (5,2.5) -- (8,2.5);
  \node[above, vertpro] at (6.5,2.5) {$d_2 = 25$ mm};
  \node[right] at (8,2.5) {Branche 1};
  % Branche basse
  \draw[thick, vertpro] (5,1.5) -- (5,0.5) -- (8,0.5);
  \node[below, vertpro] at (6.5,0.5) {$d_2 = 25$ mm};
  \node[right] at (8,0.5) {Branche 2};
  % Flèches débit
  \draw[->, thick] (2.5,1.8) -- (3.5,1.8);
  \draw[->, thick] (5.5,2.8) -- (6.5,2.8);
  \draw[->, thick] (5.5,0.2) -- (6.5,0.2);
\end{tikzpicture}
\end{center}

\textbf{Q1} \comp{APP} — Convertir le débit volumique $Q_v = 1{,}8$ m³/h en m³/s. \pts{1 pt}

\begin{reponsebox}[2.5cm]
\dotfill

\dotfill
\end{reponsebox}

\textbf{Q2} \comp{REA} — Calculer la section $S_1$ de la canalisation principale ($d_1 = 40$ mm). \pts{1 pt}

\begin{reponsebox}[3cm]
\dotfill

\dotfill

\dotfill
\end{reponsebox}

\textbf{Q3} \comp{REA} — Calculer la vitesse $v_1$ de l'eau dans la canalisation principale. \pts{1 pt}

\begin{reponsebox}[3cm]
\dotfill

\dotfill

\dotfill
\end{reponsebox}

\textbf{Q4} \comp{ANA} — Le débit se partage en deux branches identiques. Calculer le débit dans chaque branche, puis la vitesse $v_2$ dans une branche ($d_2 = 25$ mm). \pts{1 pt}

\begin{reponsebox}[4cm]
\dotfill

\dotfill

\dotfill

\dotfill

\dotfill
\end{reponsebox}

\textbf{Q5} \comp{VAL} — Les vitesses $v_1$ et $v_2$ respectent-elles la règle des 0,5 à 1,5 m/s ? Si non, proposer une solution. \pts{1 pt}

\begin{reponsebox}[3cm]
\dotfill

\dotfill

\dotfill
\end{reponsebox}

\newpage

% ================================================================
\section*{Exercice 2 — Dimensionnement électrique d'une PAC \pts{5 points}}
% ================================================================

\begin{contextebox}
Un technicien installe une pompe à chaleur (PAC) air/eau pour chauffer un immeuble. La PAC est alimentée en monophasé 230 V / 50 Hz. D'après la plaque signalétique :
\begin{itemize}[nosep]
  \item Puissance active : $P = 3\,200$ W
  \item Facteur de puissance : $\cos\varphi = 0{,}85$
\end{itemize}
La PAC fonctionne 10 heures par jour pendant 180 jours (saison de chauffe). Le prix du kWh est 0,20 €.
\end{contextebox}

% Schéma triangle des puissances
\begin{center}
\begin{tikzpicture}[scale=0.6, every node/.style={font=\small}]
  \draw[thick, bleupro] (0,0) -- (7,0) node[midway, below=4pt] {$P$ (W) — active};
  \draw[thick, rouge] (7,0) -- (7,3.5) node[midway, right=4pt] {$Q$ (VAR)};
  \draw[thick, vertpro] (0,0) -- (7,3.5) node[midway, above=4pt, sloped] {$S$ (VA) — apparente};
  % Angle
  \draw (1.5,0) arc[start angle=0, end angle=26.6, radius=1.5];
  \node at (2,0.5) {$\varphi$};
\end{tikzpicture}
\end{center}

\textbf{Q1} \comp{REA} — Calculer la puissance apparente $S$ de la PAC. \pts{1 pt}

\begin{reponsebox}[3cm]
\dotfill

\dotfill

\dotfill
\end{reponsebox}

\textbf{Q2} \comp{REA} — Calculer l'intensité du courant $I$ absorbée par la PAC. \pts{1 pt}

\begin{reponsebox}[3cm]
\dotfill

\dotfill

\dotfill
\end{reponsebox}

\textbf{Q3} \comp{ANA} — Le disjoncteur de la ligne est calibré à 20 A. Est-il correctement dimensionné ? Justifier. \pts{1 pt}

\begin{reponsebox}[2.5cm]
\dotfill

\dotfill
\end{reponsebox}

\textbf{Q4} \comp{REA} — Calculer l'énergie consommée par la PAC sur une saison de chauffe (en kWh). En déduire le coût annuel. \pts{1 pt}

\begin{reponsebox}[3.5cm]
\dotfill

\dotfill

\dotfill

\dotfill
\end{reponsebox}

\textbf{Q5} \comp{COM} — Le propriétaire demande pourquoi sa facture est élevée. Le technicien remarque que le $\cos\varphi$ est de 0,85 au lieu de 0,95 (valeur idéale). Expliquer en quoi un mauvais facteur de puissance augmente l'intensité appelée, et pourquoi c'est un problème. \pts{1 pt}

\begin{reponsebox}[4cm]
\dotfill

\dotfill

\dotfill

\dotfill

\dotfill
\end{reponsebox}

\vfill

\begin{center}
\begin{tabular}{|c|c|c|c|}
\hline
\textbf{Exercice} & \textbf{Thème} & \textbf{Questions} & \textbf{Points} \\
\hline
1 & Débit de fluide (réseau chauffage) & Q1 à Q5 & /5 \\
\hline
2 & Puissance en régime sinusoïdal (PAC) & Q1 à Q5 & /5 \\
\hline
\multicolumn{3}{|r|}{\textbf{Total}} & \textbf{/10} \\
\hline
\end{tabular}
\end{center}

\end{document}
